% !TeX root = ../main.tex

\chapter{非加性最短路径问题及其算法}

本章系统阐述非加性最短路径问题的理论基础和求解方法.从单源最短路径问题的基本概念和算法出发，引入非加性最短路径问题的数学定义；介绍线图概念和相关理论，为问题转化提供理论基础；阐明CUDA并行框架的基本原理，为后续算法的GPU实现奠定基础.

\section{预备知识}


\subsection{图的定义}

在图论中，图的定义主要由项点和边构成，一般表示为 $G(V, E, W)$，其中 $V$ 是项点集，$E$ 是所有边的集合，$W$ 是所有边权值的集合.图的总项点数一般用 $|V|=n$ 来表示，总边数则用 $|E|=m$ 表示.为了更清晰地阐述图的相关定义，引入以下网络图进行说明，具体如图2-1所示.

在图论中，根据边的指向性，可以将图分为有向图和无向图.图2-1展示了一个简单的有向图示例，边 $(\mathbf{A}, \mathbf{B})$ 表明是从点 $A$ 到点 $B$，如果边 $(\mathbf{A}, \mathbf{B})$ 没有方向指示，则为无向图.在有向图中，项点的出度表示该项点发出的边的总数量，项点入度表示该项点接收的边的总数量.在无向图中，由于没有方向指示，按照定义项点度数为与该项点相连的边的总数.

如果图中任意一条边都附加权重，则该图为有权图，否则该图为无权图.这里的权重可以是距离、时间、耗油量等.通常有权图用 $G(V, E, W)$ 表示，无权图则用 $G(V, E)$ 表示.

\subsection{图的储存格式}

图的储存格式有多种，常见的有邻接矩阵和邻接表.邻接矩阵是 \( |V| \times |V| \) 的矩阵，其中 \( A[i][j] \) 表示从节点 \( i \) 到节点 \( j \) 是否有边.邻接表是每个节点的边列表，存储了从每个节点出发的所有边.常见的网络图包括有向无环图（DAG）、带权有向图和无向图，它们在不同的应用场景中各有优势.

\section{最短路径问题}

最短路径问题是图论和算法设计中的经典问题，广泛应用于网络路由、交通规划、社交网络分析等领域.本节从最基本的单源最短路径问题出发，介绍其数学定义、关键概念及松弛操作原理，随后阐述几种经典的求解算法.这些基础知识为后续非加性最短路径问题的研究奠定理论基础.

\subsection{单源最短路径问题的定义}

\begin{definition}
在图 $G(V, E, W)$ 中，如果边 $(v_i, v_{i+1}), (v_{i+1}, v_{i+2}), (v_{i+2}, v_{i+3}), \ldots, (v_{l-1}, v_j)$ 存在，则 $l = \langle v_i, v_{i+1}, v_{i+2}, \ldots, v_j \rangle$ 就是点 $v_i$ 到点 $v_j$ 的一条完整路径.其中路径 $l$ 的距离为
\[
w(l) = \sum_{i=1}^{l-1} w(v_i, v_{i+1}).
\]
\end{definition}

\begin{definition}
最短路径是指从起点 $s$ 到点 $v_j$ 的所有路径中，使得路径长所有边权重之和最小的路径.如果存在多条最短路径，则可以选择其中任意一条作为最短路径.最短路径长度计算公式如下：
\[
d(v_i, v_j) = \min(w(L)),
\]
其中，$L$ 是从点 $v_i$ 到点 $v_j$ 的所有路径的集合，$w(L)$ 是路径集合 $L$ 中所有路径的权重的集合.
\end{definition}

\begin{definition}
设 $u$ 和 $v$ 是图 $G(V, E, W)$ 中两个顶点，$w(u,v)$ 表示顶点 $u$ 到点 $v$ 的边的权重，$d(u)$ 和 $d(v)$ 分别表示源点 $s$ 到顶点 $u$ 和顶点 $v$ 的最短距离.如果 $d(u) + w(u,v) < d(v)$，则更新 $d(v)$ 为 $d(u) + w(u,v)$.这一过程就是松弛操作.
\end{definition}

单源最短路径问题是求解从源点到图中所有其他顶点的最短路径，求解也称为松弛操作.

\subsection{经典最短路径问题算法}

经典的最短路径问题算法主要包括Dijkstra算法、Bellman-Ford算法和Floyd-Warshall算法.其中，Dijkstra算法因其高效的性能和广泛的适用性，成为解决单源最短路径问题的首选方法.

\subsubsection{Dijkstra 算法}

Dijkstra算法由荷兰计算机科学家E.W. Dijkstra于1956年提出，是解决单源最短路径问题最常用的贪心算法.该算法适用于所有边权重均为非负的有向图或无向图.

**算法思想：** Dijkstra算法采用贪心策略，每次从未访问的顶点中选择距离源点最近的顶点，将其加入已访问集合，并据此更新其他顶点的最短距离.通过对所有顶点的这样处理，最终可得到从源点到所有其他顶点的最短路径.

**时间复杂度：** 采用二叉堆实现的Dijkstra算法时间复杂度为 $O((|V| + |E|) \log |V|)$；采用线性查找实现的Dijkstra算法时间复杂度为 $O(|V|^2)$.

**算法伪代码：**

\begin{algorithm}
  \SetAlgoLined
  \caption{Dijkstra Algorithm}
  \label{algo:dijkstra}
  
  \textbf{Input:} Graph $G = (V, E, W)$, source vertex $s$\\
  \textbf{Output:} Shortest distances $d(v)$ for all $v \in V$\
  
  \textbf{Initialization:}\\
  \For{each $v \in V$}{
    $d(v) \leftarrow \infty$\
  }
  $d(s) \leftarrow 0$\
  $T \leftarrow V/s$\
  
  \While{$T$ is not empty}{
    \textbf{Min:} $u \leftarrow \{v : v \in T \text{ and } \forall c \in T, d(v) \leq d(c)\}$\
    
    \If{$d(u) = \infty$}{
      \textbf{break}\
    }
    
    Remove $u$ from $T$\
    
    \textbf{Relax:}\
    \For{each $(u, v) \in E$}{
      \If{$d(u) + w(u, v) < d(v)$}{
        $d(v) \leftarrow d(u) + w(u, v)$\
      }
    }
  }
  
  \Return{$d$}\
\end{algorithm}

在上述伪代码中，$d(v)$ 表示源点 $s$ 到顶点 $v$ 的最短距离，$T$ 表示未处理顶点的集合，$w(u, v)$ 表示边 $(u, v)$ 的权重.

**算法步骤说明：** 
\begin{enumerate}
  \item **初始化**：将所有顶点的距离初始化为无穷大，源点距离为0，未处理顶点集合 $T$ 为除源点外的所有顶点.
  \item **主循环**：当存在未处理的顶点时，重复以下步骤：
  \begin{enumerate}
    \item 从未处理顶点集合 $T$ 中选择距离最小的顶点 $u$.
    \item 若 $d(u) = \infty$，说明剩余顶点不可达，退出循环.
    \item 将顶点 $u$ 从集合 $T$ 中移除.
    \item 对从 $u$ 出发的所有边执行松弛操作，若发现更短路径则更新距离.
  \end{enumerate}
  \item **输出**：返回所有顶点的最短距离数组 $d$.
\end{enumerate}

\subsubsection{Bellman-Ford 算法}

**算法特点：** Bellman-Ford算法相比Dijkstra算法的优势在于能够处理边权重为负的图，并可以检测图中是否存在负权环.其缺点是时间复杂度较高，为 $O(|V| \cdot |E|)$.该算法通过对所有边进行多次松弛操作来计算最短路径.

\subsubsection{Floyd-Warshall 算法}

**算法特点：** Floyd-Warshall算法用于求解所有顶点对之间的最短路径，采用动态规划的思想.时间复杂度为 $O(|V|^3)$，空间复杂度为 $O(|V|^2)$.该算法的优于点包括算法实现简单，且能处理负权重的边（但不能处理负权环）



\section{非加性最短路径问题}

非加性最短路径问题是传统最短路径问题的推广和扩展，其路径权重的定义不满足加性性质.本节从非加性的基本定义出发，介绍问题的分类方法，进而阐述相应的求解算法.

\subsection{非加性的定义}

线图是图论中的一个重要概念，它为将非加性最短路径问题转化为标准问题提供了理论基础.通过构造适当的线图，可以将某些非加性问题转换为加性问题，从而利用现有的算法进行求解.

给定一个图 \( G = (V, E) \)，其线图 \( L(G) = (V', E') \) 是一个图，其中顶点集 \( V' = E \) 是原图的边集合，边集 \( E' \) 中的两个顶点（即原图的两条边）相邻当且仅当它们在原图中共享一个公共顶点.线图的构造允许我们从不同的视角重新审视原图的结构，在某些情况下可以简化问题的求解.

线图的性质对于问题转化至关重要.例如，原图中相邻的顶点对应线图中相邻的路径，这种对应关系使得我们可以在线图上进行最短路径计算，其结果可以映射回原图中的非加性最短路径问题.

\section{CUDA并行框架介绍}

CUDA并行框架是NVIDIA开发的通用并行计算平台和API，它允许开发者在GPU上执行大规模并行计算任务.CUDA框架的基本概念和编程模型为后续第三章的算法GPU并行实现提供了技术基础.

CUDA编程模型主要基于以下核心概念：
\begin{itemize}
  \item **线程（Thread）**：线程是CUDA并行计算中最基本的执行单位，上万甚至数百万个线程可以并发执行.
  \item **线程块（Block）**：多个线程组成一个块，同一块内的线程可以协作共享内存和进行同步.
  \item **网格（Grid）**：多个线程块组成一个网格，网格是GPU执行的最高层次组织.
  \item **GPU全局内存**：所有块和线程都可以访问的内存，适合存储大规模数据.
  \item **共享内存**：块内线程共享的高速缓存，用于线程间通信和同步.
\end{itemize}

在最短路径问题的并行计算中，CUDA框架的这些特性可以被充分利用.例如，在计算多源最短路径时，可以为每个源点分配一个线程块，在块内进行并行化的Dijkstra或Bellman-Ford算法迭代.这样可以显著加速问题求解，特别是对于大规模图的处理.

\section{本章小结}

本章系统阐述了非加性最短路径问题的理论基础和求解方法.首先介绍了图论的基本概念，包括图的定义和储存格式；其次阐述了传统的单源最短路径问题及其经典求解算法；再次引入了非加性最短路径问题的概念，并根据路径权重函数的性质进行了分类，阐述了不同类型问题的求解方法；最后介绍了线图的基础理论和CUDA并行框架，为后续的算法设计和实现奠定了基础.

与传统最短路径问题相比，非加性最短路径问题的主要挑战在于路径权重不满足加性性质，这使得传统的松弛操作失效.通过引入线图、修改松弛条件或采用动态规划等方法，可以在一定程度上克服这些困难.同时，利用CUDA等并行计算框架，可以进一步提高算法的计算效率，特别是对于大规模图的处理.

在后续的研究中，我们将在这些基础理论和技术的基础上，设计针对具体非加性最短路径问题的高效算法，并通过GPU并行实现来进一步提升计算性能.
